\begin{table}[htbp]\centering\small
\caption{最终分问配置的334天开发期闭环评价结果}
\label{tab:annual_summary}
\begin{tabularx}{\linewidth}{LRRRR}\toprule
指标 & 问题2 & 问题3 & 问题4日前版 & 问题4滚动版\\\midrule
计划费用/万元 & 1350.80 & 1350.31 & 1428.01 & 1426.52\\
调整费用/万元 & 0.00 & 24.87 & 0.00 & 28.31\\
紧急费用/万元 & 50.48 & 11.97 & 51.40 & 12.19\\
实际总费用/万元 & 1401.28 & 1387.15 & 1479.41 & 1467.02\\
紧急电量/kWh & 85017.69 & 21029.90 & 80544.00 & 21293.41\\
紧急购电天数 & 113 & 119 & 106 & 116\\
紧急购电连续时段数 & 205 & 204 & 195 & 193\\
日总费用CVaR95/元 & 70731.74 & 65552.76 & 83044.31 & 77575.87\\
最高日费用/元 & 89306.82 & 67941.61 & 108080.37 & 81797.18\\
2月1日期初库存/kWh & 8875.19 & 8814.60 & 9251.54 & 9030.40\\
12月31日期末库存/kWh & 6240.32 & 6240.32 & 6269.86 & 6274.99\\
\bottomrule\end{tabularx}
\par\smallskip\parbox{\linewidth}{\footnotesize 问题2与问题4日前版沿用母模型；问题3与问题4滚动版在零时及日内均启用候选筛选。预测器与实时执行器沿用原冻结配置，筛选适用问题与时点依据开发期比较确定。实际费用不含候选评分的终端偏离项；CVaR95为实际日总账单的经验尾部均值。}
\end{table}
